\documentclass[../00_Main.tex]{subfiles}
\begin{document}
\appendix
\section{Wavelength Calculation at 100 Hz}
\label{appendix:wavelength}

To estimate the electromagnetic wavelength $\lambda$ at 100 Hz, we start with the general wave relation:

\[
\lambda = \frac{v}{f}
\]

In a medium, the wave speed is reduced due to the material's permittivity \( \varepsilon \) and permeability \( \mu \). This gives:

\nota{source Griffiths,Introduction to Electrodynamics, Sec. 9.2.2 \url{https://nucleares.unam.mx/~martinel/griffiths_4ed.pdf} } 

\[
v = \frac{c}{\sqrt{\varepsilon_r \mu_r}}
\quad \Rightarrow \quad
\lambda = \frac{c}{f \sqrt{\varepsilon_r \mu_r}}
\]

Where:
\begin{itemize}
  \item $c = 3 \times 10^8$ m/s is the speed of light in vacuum
  \item $f$ is frequency in Hz
  \item $\varepsilon_r$ is the relative permittivity
  \item $\mu_r$ is the relative permeability (typically $\mu_r \approx 1$ in non-magnetic media)
\end{itemize}

\subsection{Wavelength in Gelatin–Salt–Water Phantom}

We note that the permittivity of saltwater is lower than that of pure water, as dissolved ions disrupt the orientational polarization of water molecules.\nota{\url{https://doi.org/10.1063/1.1746645} and  \url{https://arxiv.org/pdf/1208.5169}}

Using the wavelength relation:

\[
\lambda = \frac{c}{f \sqrt{\varepsilon_r \mu_r}}
\]

it follows that if $\varepsilon_{\text{saltwater}} < \varepsilon_{\text{water}}$, then:

\[
\sqrt{\varepsilon_{\text{saltwater}}} < \sqrt{\varepsilon_{\text{water}}} \quad \Rightarrow \quad \lambda_{\text{saltwater}} > \lambda_{\text{water}}
\]

In other words, the wavelength in saltwater will be longer than that in pure water, all else equal.

To provide a conservative estimate, we compute the wavelength in pure water using:

\begin{itemize}
  \item $f = 100$ Hz
  \item $\varepsilon_r \approx 78.5$ at 25°C \nota{\url{https://nvlpubs.nist.gov/nistpubs/jres/56/jresv56n1p1_a1b.pdf} - Dielectric Constant of Water from 0 0 to 1000 C
C. G. Malmberg and A. A. Maryott}
  \item $\mu_r \approx 1$
\end{itemize}

\[
\lambda_{\text{water}} \approx \frac{3 \times 10^8}{100 \cdot \sqrt{78.5}}~\text{m}
= 338 599~\text{m} = 338.6~\text{km}
\]

Thus, the wavelength in the gelatin–salt–water phantom is expected to be at least 338.6 km, and likely longer — reinforcing the validity of the quasi-static approximation for the frequencies and scales of interest.

\vspace{0.5em}
\subsection{Wavelength in Protopasta C-Doped PLA}

For Protopasta conductive PLA, a carbon-doped 3D printing filament:
\begin{itemize}
  \item $f = 100$ Hz
  \item $\varepsilon_r \sim 10$ \url{https://doi.org/10.1016/j.compositesb.2016.04.030}
  \item $\mu_r \approx 1$ (we take the approximation for non-magnetic materials).
\end{itemize}

\[
\lambda_{\text{PLA}} = \frac{3 \times 10^8}{100 \cdot \sqrt{10}} ~\text{m} \approx 948683 ~\text{m} = 948.6~\text{km}
\]

\vspace{0.5em}
\noindent
In both materials, the resulting wavelengths are several orders of magnitude larger than the physical system size. This justifies the quasi-static approximation used in Section~\ref{subsec:pde} for all EEG-relevant frequencies.


\biblio % Needed for referencing to working when compiling individual subfiles - Do not remove
\end{document}